\tutorial{4}{Coupling an existing MODFLOW 6 model}

\begin{goal}
A colleague has built a MODFLOW~6 model of Tutorial Creek: in feet, on a grid of its own, with its own recharge, ET,
rivers, well and head boundary. You couple the Raven model of Tutorial~2 to it without changing it: link the HRUs to its
cells, let Raven take over recharge and ET, send its river flows to Raven's reaches, then add water-table ET from Raven, a
wetland and the reservoir, and check the result. Time: about 40 minutes.
\end{goal}

\section{The two models}
\begin{shell}
$ cp -r $TUT/t4_existing_model/start ~/work/t4
$ cd ~/work/t4
\end{shell}
The folder holds the Raven model of Tutorial Creek as it was at the start of Tutorial~2 (no groundwater), the HRU polygons
\cmd{hrus.geojson}, and the MODFLOW~6 model in \cmd{model/}:
\term{t4_ls_model.txt}
\term{t4_nam.txt}
\begin{itemize}
\item A regular grid (\cmd{DIS}) of \V{t4:nrow}~rows and \V{t4:ncol}~columns of \V{t4:cellft}-ft cells
      (\V{t4:cellm}~m), two layers, lengths in feet; its origin is given in feet in UTM zone 11N:
\term{t4_dis_head.txt}
\item A one-day steady-state period followed by monthly periods, \V{t4:nper} in all, from 1 January 2010:
\term{t4_tdis_head.txt}
\item Its own recharge (\cmd{RCH\_MODEL}) and evapotranspiration (\cmd{EVT\_MODEL}), as arrays.
\item \V{t4:nriv} river cells (\cmd{RIV\_CREEK}) whose auxiliary variable \cmd{SUBBASIN} names the Raven subbasin of each reach:
\term{t4_riv_head.txt}
\item The pumping well (\cmd{WEL\_PUMP}, monthly rates) and a head boundary at the outlet (\cmd{GHB\_OUTLET}, \V{t4:nghb} cells).
\end{itemize}

\section{Link the Raven model to it}
\subsection{The Raven files}
The \cmd{.rvh} and \cmd{.rvi} edits are those of Tutorial~2: mark the coupled HRUs with an aquifer profile and send their
percolation to \cmd{GROUNDWATER}.
\diff{t4_a_rvh.diff}
\diff{t4_a_rvi.diff}
The \cmd{.rvp} needs no aquifer classes: the layers and their properties are the model's. A profile is written as
\cmd{name, 0}; it only marks the coupled HRUs and carries the exchange settings:
\diff{t4_a_rvp.diff}

\subsection{The coupling file}
Create \cmd{Creek.rvg}:
\lstinputlisting[style=file]{generated/t4_rvg_a.txt}
\begin{itemize}
\item \cmd{:MF6Simulation} names the simulation file. Raven copies the folder that holds it into \cmd{out/mf6/} and runs the
      copy; \cmd{model/} is never changed.
\item \cmd{:MF6CRS} says in which projection the model's coordinates are. Raven reads the cell outlines from the model
      (\cmd{DELR}, \cmd{DELC}, \cmd{XORIGIN}, \cmd{YORIGIN}, \cmd{ANGROT}), converts feet to metres, and places them on the map
      under the HRU polygons. (Without a projection, \cmd{:MF6GridFile} --- cell polygons --- or \cmd{:OverlapWeights} ---
      a table of HRU/cell overlaps --- link them instead.)
\item \cmd{:MF6StartDate} is the calendar date of the model's time zero: Raven keeps the model's stress periods and splits
      each into Raven days.
\item \cmd{:RavenTakesOver RCH\_MODEL EVT\_MODEL}: every recharge, ET or UZF package must be declared taken over (removed:
      Raven supplies recharge) or kept (\cmd{:KeepPackages}); an undeclared one stops the run, so recharge is never counted
      twice.
\item \cmd{:StreamPackage RIV\_CREEK AUX SUBBASIN}: the river package whose flows go to Raven's reaches, with the subbasin of
      each cell from its auxiliary variable.
\item \cmd{:SeepageToRaven ON} adds a drain (\cmd{DRN\_RAVEN}) at the top of the coupled columns; what it discharges returns
      to Raven.
\end{itemize}

\subsection{Run}
\term{t4_run_a.txt}
\term{t4_summary.txt}
\begin{tocheck}
\begin{itemize}
\item \cmd{existing model}: the model, its grid and its units (\emph{feet, days}).
\item \cmd{packages taken over} and \cmd{packages in balance} (the model's own packages that stay: their net flow is the
      column \cmd{other model packages} of \cmd{GWBudget.csv}).
\item \cmd{HRU coverage}: \V{t4:coverage}\% of the coupled HRUs' area lies on active cells. The rest lies outside the model:
      its edges do not follow the HRU outlines and cut off the outer parts of the slope HRUs. A low value usually means a
      wrong projection or wrong length units.
\item The second part of the file lists each HRU with \cmd{linked\_fraction\_of\_polygon}: the slope HRUs 3, 7 and 8
      are linked at \V{t4:link3}, \V{t4:link7} and \V{t4:link8}\%, the others at \V{t4:linkothers}\% or more.
\end{itemize}
\end{tocheck}
\begin{trouble}
\begin{itemize}
\item \emph{``\dots must be named in :RavenTakesOver or :KeepPackages''}: declare every recharge, ET and UZF package.
\item \emph{``a stream cell whose subbasin does not exist''}: the \cmd{SUBBASIN} values must be subbasin IDs of the \cmd{.rvh}.
\item \emph{``each stress period must contain a whole number of Raven time steps''} or \emph{``\dots longer than the model''}:
      check \cmd{:MF6StartDate}, \cmd{:StartDate} and \cmd{:Duration}.
\item Grid options of the generated mode (\cmd{:GridCellSize}, \cmd{:LandSurface}, \cmd{:Wells}, \dots) in this \cmd{.rvg}
      stop the run: the model has its own.
\end{itemize}
\end{trouble}
The working copy shows what Raven did to the model's name file:
\term{t4_mf6nam.txt}

\subsection{Units}
The model works in feet and days; Raven converts at the interface (1~ft = 0.3048~m, 1~ft\textsuperscript{3} =
0.0283168~m\textsuperscript{3}). Its recharge, seepage, river and package flows are written to \cmd{GWBudget.csv} in
m\textsuperscript{3}/d and heads in metres, as for a generated model. A model in seconds or in centimetres is converted the
same way; a model that does not say its units is taken as metres and days.

\subsection{Check the balance}
\term{t4_audit_a.txt}
\begin{tocheck}
The recharge identity is 0 (every cubic metre Raven sent reached the model's cells, in feet), the cumulative residual
\V{t4:a:auditRes}, no unconverged steps (\V{t4:a:auditUnconv}).
\end{tocheck}
Mean flows (m\textsuperscript{3}/d): recharge \V{t4:a:rech}, seepage to Raven \V{t4:a:seep}, discharge to the creek
\V{t4:a:toriv}, leakage from the creek \V{t4:a:fromriv}, the model's own well and boundary \V{t4:a:other}, storage release
\V{t4:a:stor}.
\begin{note}[Look at the storage]
The storage release is large and positive: \V{t4:a:stor1}~m\textsuperscript{3}/d on average in the first year, still
\V{t4:a:stor3} in the third. The model's first period is a steady state, solved with the recharge of Raven's first day ---
\V{t4:a:rechday1}~m\textsuperscript{3}/d (\V{t4:a:rechday1mm}~mm over the coupled HRUs), because the initial soil moisture
of the \cmd{.rvc} lets the subsoil drain fast. The model starts with too much water and drains for years. In a real
study start Raven from spun-up soil states (a \cmd{solution.rvc} of an earlier run), or let the colleague's steady period
see a representative recharge.
\end{note}

\section{Water-table ET from Raven}
The model's own ET was removed with \cmd{EVT\_MODEL}. Raven can take ET from the water table itself: an
\cmd{EXTINCTION\_DEPTH} in the profiles adds \cmd{EVT\_RAVEN} at the top active cells; its rate is Raven's unmet
evaporative demand of the HRU, falling to zero at the extinction depth below the cell top.
\diff{t4_b_rvp.diff}
\term{t4_run_b.txt}
\term{t4_summary_b.txt}
Water-table ET takes \V{t4:b:et}~m\textsuperscript{3}/d on average (up to \V{t4:b:etmax} in summer), mostly where the
water table is shallow; seepage to Raven falls to \V{t4:b:seep}~m\textsuperscript{3}/d.
\begin{trouble}
\emph{``EXTINCTION\_DEPTH > 0 \dots the model keeps its EVT \dots package''}: ET from the model and from Raven would both act;
take the model's EVT or UZF package over, or set \cmd{EXTINCTION\_DEPTH} to 0.
\end{trouble}

\section{The wetland and the reservoir}
Tutorial Creek's wetland (HRU~9) stores water in Raven's \cmd{DEPRESSION} store, and CreekLake (HRU~5) is a reservoir.
Both can exchange water with the aquifer under them:
\begin{itemize}
\item \cmd{:SurfaceWaterExchange DEPRESSION WetlandHRUs LEAKANCE 0.01} adds \cmd{SWX\_RAVEN}, a river-type boundary at the
      cells under the HRU group \cmd{WetlandHRUs}: its stage is the water level of the depression store, its conductance
      the leakance [1/d] times the area, and it never takes more water than the store holds.
\item \cmd{:ReservoirExchange} couples every reservoir that has \cmd{:SeepageParameters} and a lake HRU (\cmd{:HRUID}) with an
      aquifer profile: the water table under the lake becomes the reservoir's local groundwater head, and the reservoir's
      seepage enters the aquifer the next day (\cmd{RES\_RAVEN}). The stages must be absolute (\cmd{:AbsoluteCrestHeight}).
\end{itemize}
Define the group and give the reservoir its seepage parameters (seepage coefficient and reference head):
\diff{t4_c_rvh.diff}
\diff{t4_c_rvi.diff}
\diff{t4_c_rvg.diff}
\term{t4_run_c.txt}
\term{t4_summary_c.txt}
\term{t4_mf6nam_c.txt}
All \V{t4:nraven} packages Raven adds are in the working copy, each marked as added by Raven. Mean flows (m\textsuperscript{3}/d):
the wetland leaks \V{t4:c:swin} into the aquifer (and receives \V{t4:c:swout}); CreekLake sends \V{t4:c:ressent}
of seepage, all of it delivered (\V{t4:c:resdel}); water-table ET \V{t4:c:et}; seepage to Raven \V{t4:c:seep}.

\subsection{Final checks}
\term{t4_audit_c.txt}
\begin{tocheck}
The reservoir lines: \cmd{delivered / sent} 1 and no seepage left undelivered; the surface-store identity 0; recharge
identity 0; cumulative residual \V{t4:c:auditRes}; \V{t4:c:auditUnconv} unconverged steps.
\end{tocheck}
\term{t4_plot.txt}
\fig{t4_gw_budget.png}{The budget of the coupled existing model: the model's own well and boundary (\emph{model's own
packages}), water-table ET and the storage release of the drain-down from its steady first day.}
\fig[0.72]{t4_gw_map.png}{Water-table depth below the model's top on the model's own cells (the depths are converted from feet).}

\section{Summary}
Coupling an existing model takes the \cmd{.rvh} and \cmd{.rvi} edits of any coupled model, profiles without layers in the
\cmd{.rvp}, and a short \cmd{.rvg}. Raven never changes your model: \cmd{out/mf6/} is a working copy with the stress periods
split into days, the taken-over packages commented out and its own packages added. The finished state is
\cmd{\$TUT/t4\_existing\_model/finished}. The manual's chapter ``Coupling an existing MODFLOW~6 model'' describes restarts
and later starts, DISV grids, the water mover, steady-state models and the safeguards; the package's
\cmd{examples/Liard\_existing\_model} is a larger example.
